\begin{textsample}{NEG Dim 1 – vem\_ed\_02 – Score 20.00 – vem\_ed\_02/t008.txt}  \label{ex:f1_neg_029}
SAÚDE E TECNOLOGIA Iniciado no segundo semestre de 2019, o PCI"Qualidade em \textbf{equipamentos} \textbf{médico-hospitalares}: políticas e \textbf{procedimentos}"envolve alunos de Sistemas Biomédicos da Fatec Bauru e do Programa de Gestão em Saúde da instituição norte-americana BridgeValley Community \& Technical College.
Os professores brasileiros responsáveis pelo projeto são Camila Maria da Costa Kami, de inglês, e Luiz Roberto Madureira Iório, do curso de Sistemas Biomédicos.
Nos Estados Unidos, Liesa Kyer, professora da área de Enfermagem, conduz as atividades.
Em 2019, os grupos de estudantes (10 norte-americanos e 17 brasileiros) mencionaram os \textbf{equipamentos} que conheciam e compararam marcas e preços nos dois países.
Cada equipe escolheu um \textbf{equipamento} e levantou informações sobre seu funcionamento, \textbf{operação} e manutenção.
A segunda edição do projeto, neste segundo semestre de 2020, \textbf{contemplará} o contexto da COVID-19: os alunos irão discutir as situações vivenciadas nos dois países e, principalmente, os protocolos de higienização dos \textbf{equipamentos} \textbf{médicos}."Os PCIs propostos pelo Centro Paula Souza abriram a possibilidade de internacionalização de projetos, como uma nova e importante perspectiva para alunos e docentes da Fatec Bauru, pois permitiram que diferentes culturas e estruturas curriculares ' conversassem ' entre si, trocando experiências e conhecimentos.
O impacto tem sido e continuará sendo extremamente positivo para a unidade e gera mais expectativas para que todos os cursos possam \textbf{contemplar} mais projetos em suas áreas", ressalta o diretor da Fatec Bauru, Sebastião Gândara Vieira. ' Os alunos tiveram a oportunidade de praticar a língua inglesa e conhecer um pouco sobre a cultura do parceiro; foi uma ótima oportunidade para desenvolverem a competência intercultural, habilidade indispensável nos dias atuais ', afirma Camila Kami.
Luiz Roberto Madureira Iório completa:"A oportunidade de trabalhar em equipe com a professora \textbf{norte-americana} ajudou a estreitar laços entre Fatec Bauru e BridgeValley Community \& Technical College".
Liesa Kyer está animada com o prosseguimento dos trabalhos."Meus alunos aprenderam a colaborar com os brasileiros em um projeto que demanda pesquisa e tempo; para a maioria deles, esta é uma oportunidade única na vida".
CORREÇÃO Na edição número 1 de VEm com PCI, houve um problema na editoração eletrônica do texto"Parceiros alinhados e comprometidos", e o trecho com o depoimento de Adriane Monteiro Fontana, diretora da Fatec São Caetano do Sul, ficou truncado.
Sobre o PCI"Education and Leadership", desenvolvido no primeiro semestre de 2020 pela Cesu, entre professores, coordenadores e diretores de Fatecs e estudantes de pós-graduação da Florida International University (FIU), a diretora Adriane comentou:"Excelente oportunidade de intercâmbio de ideias e boas práticas entre as Fatecs e instituições de ensino superior em outros países; nesse sentido, é possível conversar com pares para que tal contribuição tenha um significado efetivo".

% matched lemmas: contemplar, equipamento, médico, médico-hospitalar, norte-americana, operação, procedimento
\end{textsample}
